\documentclass[../main.tex]{subfiles}

\begin{document}

\section{項・述語関数・命題}

本章では，論理学において文をどのように分析するかを確認する。文はまず主語と述語をもつものとして扱われ，主語は対象を指し，述語はその対象について成り立つ性質や関係を表す。この見方により，日本語の文も形式的な論理式へ写し取ることができる。

\subsection{主語と述語}

最も基本的な文の形は，主語 \(S\) と述語 \(P\) によって
\[
  \text{すべての } S \text{ is } P
\]
または
\[
  \text{ある } S \text{ is } P
\]
と表される。ここで \(S\) は主語，\(P\) は述語である。述語は対象について「\(P\) である」または「\(P\) でない」と述べる部分であり，文の真偽を決める働きをもつ。

\begin{sidebox}[TBred]{注意}
  日本語の文が本当にすべて主語と述語でできているのか，という点は検討を要する。日常言語には主語が省略される文や，文脈によって主語が補われる文が多いためである。ただし，論理学では分析のために主語と述語の形へ整える。
\end{sidebox}

\subsection{個体・性質・関係}

論理学では，対象を指す表現を項として扱う。項が指すものを個体という。固有名詞は典型的な項であり，個々の対象を指示する。
\[
  \text{個体} \ (\mathrm{individuality})
  \quad \longleftarrow \quad
  \text{term, 項}
  \quad
  (\text{固有名詞})
\]
このような項を関数的に記述することが，述語論理の基本的な発想である。

性質は，ある個体に対して成り立つかどうかを問うことができる。個体 \(x\) が性質 \(\varphi\) をもつことを
\[
  \varphi(x)
  =
  x \text{ は } \varphi \text{ である}
\]
と表す。この \(\varphi\) は単項述語関数である。

これに対して，複数の個体のあいだに成り立つものを関係という。たとえば，個体 \(x,y\) のあいだに関係 \(\psi\) が成り立つことを
\[
  \psi(x,y)
  =
  x \text{ と } y \text{ とは } \psi \text{ という関係にある}
\]
と表す。一般に，\(n\) 個の対象のあいだの関係は
\[
  \zeta(x_1,\ldots,x_n)
\]
と表され，これは \(n\) 項述語関数である。

\subsection{命題と内包・外延}

1つの物体や人物は，複数の固有名詞や記述によって指されることがある。たとえば，
\[
  \Delta ABC \text{ は二等辺三角形である}
\]
と
\[
  \Delta ABC \text{ は二等角三角形である}
\]
という2つの命題を考える。この2つは同じ対象について述べているが，同じ意味であるかどうかは別に検討しなければならない。
\[
  \left.
  \begin{array}{l}
    \Delta ABC \text{ は二等辺三角形である} \\
    \Delta ABC \text{ は二等角三角形である}
  \end{array}
  \right\}
  \quad
  \text{同じ意味か}
\]

ここで重要になるのが，内包と外延の区別である。内包とは対象の記述であり，外延とはその記述に該当する対象の集合である。
\[
  \begin{array}{rl}
    \text{内包} \ (\mathrm{intension})
      &= \text{対象の記述} \ (\mathrm{description}) \\[2mm]
    \text{外延} \ (\mathrm{extension})
      &= \text{対象となっている集合}
  \end{array}
\]
二等辺三角形と二等角三角形は，定義としての記述は異なるが，通常のユークリッド幾何においては同じ三角形の集合を定める。このように，内包が異なっても外延が一致する場合がある。

\section{述語関数の解釈と同一性}

\subsection{述語関数の解釈}

述語関数は，個体の組を真理値へ対応させるものとして解釈できる。たとえば
\[
  \psi(x,y,z)
\]
は，3つの個体 \(x,y,z\) のあいだに関係 \(\psi\) が成り立つかどうかを表す。したがって，その外延は
\[
  \left\{
    \langle u,v,w \rangle
    \mid
    u,v,w \in A_{\psi}
  \right\}
\]
のような順序組の集合として理解できる。

\subsection{不可識者同一性原理}

不可識者同一性原理とは，識別できないものは同一である，という原理である。より形式的には，すべての性質が一致する2つの対象は同一である，と述べられる。
\[
  \text{性質がすべて一致するものは同一である}
\]
この考え方は，任意の個体 \(x,y\) と任意の単項述語関数 \(\varphi\) について，次のように表される。
\[
  \begin{array}{c|c}
    \text{任意の } x,y \text{ について}
    &
    \text{任意の単項述語関数 } \varphi \text{ について}
    \\ \hline
    x = y \ \df
    &
    \varphi(x) = \varphi(y)
  \end{array}
\]
これは，同一性に関する分入則へつながる。すなわち，ある対象と別の対象が同一であれば，一方について成り立つ性質は他方についても成り立つ。

\subsection{述語関数の同値性分入則}

2つの述語関数 \(\varphi,\psi\) が同値であるとは，任意の個体定項 \(a_i\) について，それらを適用した命題の真理値が一致することである。
\[
  \begin{array}{c|c}
    \text{任意の単項述語関数 } \varphi,\psi \text{ について}
    &
    \text{任意の個体定項 } a_i \text{ について}
    \\ \hline
    \varphi \equiv \psi \ \df
    &
    \varphi(a_i) \equiv \psi(a_i)
  \end{array}
\]
したがって，同値な述語関数は，個体を代入したあとでも同じ真理値を与える。

\subsection{真理値をもつ表現}

真理値をもつ表現は文，すなわち命題である。たとえば
\[
  \varphi(a,b)
  =
  a \text{ と } b \text{ とは } \varphi \text{ という関係にある}
\]
は真または偽のいずれかの値をもつ。文を指す記号としては
\[
  A,B,C,\ldots
\]
を用いる。

\section{命題と真理関数的文脈}

\subsection{命題に関する同値性分入則}

命題に関する同値性分入則は，命題 \(P,Q\) が同値であるならば，任意の真理関数的文脈 \(\Sigma\) の中で \(P\) を \(Q\) に置き換えても，全体の真理値は変わらない，という規則である。
\[
  \begin{array}{c|c}
    \text{任意の命題記号 } P,Q \text{ について}
    &
    \text{任意の真理関数的文脈 } \Sigma \text{ について}
    \\ \hline
    P \equiv Q
    &
    \Sigma[P] \equiv \Sigma[Q]
  \end{array}
\]
つまり，真理値が等しければ，それらを入れ替えた文脈全体の真理値も等しくなる。

\begin{center}
\begin{tikzpicture}[>=Stealth]
  \node at (0,0) {\(\lambda p\)};
  \node[draw, rectangle, minimum width=2cm, minimum height=1cm] at (3,0) {関数};
  \node at (6,0) {\(\lambda p\)};
  \draw[->] (0.6,0) -- (2,0);
  \draw[->] (4,0) -- (5.4,0);
  \node at (0,-0.45) {\(\binom{1}{0}\)};
  \node at (6,-0.45) {\(\binom{1}{0}\)};
\end{tikzpicture}
\end{center}

\subsection{\texorpdfstring{\(n\)}{n} 種類の命題記号を含む論理的文脈}

\(n\) 種類の命題記号を含む論理的文脈は，真理値の組を入力とし，真理値を出力する関数として表される。
\[
  F
  =
  \Omega^n
  =
  \left\{
    \langle p_1,p_2,\ldots,p_n\rangle
    \mid
    p_1,p_2,\ldots,p_n \in \Omega = \{1,0\}
  \right\}
  \longrightarrow
  \Omega = \{1,0\}
\]
したがって，このような文脈は真理関数的文脈とみなされる。入力が一意的に決まれば，出力も一意的に決まらなければならない。そして，その出力は \(1\) または \(0\) の真理値である。
\[
  \therefore
  \quad
  \text{命題は真理値が入力となる}
\]

\end{document}
